% !TEX root =./ECE343_Assessment_main.tex

\begin{questions}
	
%%%%%%%%%%%%%%%%%%%%%%%%%%%%%%%%%%%%%%%%%
% 3.1
%%%%%%%%%%%%%%%%%%%%%%%%%%%%%%%%%%%%%%%%%
\fullwidth{\textbf{Learning Objective 3.1: } I can describe the origin of the transmission line equation, understand the method of developing a solution, and interpret the solution.}

%\question By using direct substitution, prove that the forward traveling portion of the voltage and current waves:
%		\eq{
%			v(z,t) & = V_0^+ e^{-\alpha z} \cosp{\omega t - \beta z} +  V_0^- e^{\alpha z} \cosp{\omega t + \beta z} \\
%			i(z,t) & = I_0^+ e^{-\alpha z} \cosp{\omega t - \beta z} +  I_0^- e^{\alpha z} \cosp{\omega t + \beta z}
%		}
%		is a solution to the Telegrapher's Equations. 
\question Develop a program/script to plot the reverse traveling voltage wave on a lossy transmission line. Choose reasonable values for $\omega$ and $\beta$. Provide a graphic for 4 different time steps across a period. Describe the behavior you observe over these time steps. 

\begin{solution}
	\scalebox{0.7}{
		\begin{tikzpicture}
			\begin{axis}[domain=0:60e-3, 
				samples=200, 
				xmin=0, 
				xmax=65e-3, 
				ymax=10,
				ymin=-10,
				title={Reverse Wave, $t = 0$},
				scaled x ticks=base 10:3,
				xtick scale label code/.code={},
				minor x tick num = 1,
				minor y tick num = 1,
				yticklabel={\empty},
				xticklabel={\empty},
				xlabel=$z$,
				ylabel={$V(z,t)$},
				axis y line = middle,
				axis x line = middle,
				x label style={xshift=0.5cm, yshift=-0.5cm},
				y label style={xshift=-0.5cm, yshift=0.75cm},
				clip=false,
				]
				\addplot[blue] {0.5*exp(50*x)*cos(360*50*x + 0)} ;
			\end{axis}
		\end{tikzpicture}
		} % end scalebox
		
		%T/4
			\scalebox{0.7}{
		\begin{tikzpicture}
			\begin{axis}[domain=0:60e-3, 
				samples=200, 
				xmin=0, 
				xmax=65e-3, 
				ymax=10,
				ymin=-10,
				title={Reverse Wave, $t = \frac{T}{4}$},
				scaled x ticks=base 10:3,
				xtick scale label code/.code={},
				minor x tick num = 1,
				minor y tick num = 1,
				yticklabel={\empty},
				xticklabel={\empty},
				xlabel=$z$,
				ylabel={$V(z,t)$},
				axis y line = middle,
				axis x line = middle,
				x label style={xshift=0.5cm, yshift=-0.5cm},
				y label style={xshift=-0.5cm, yshift=0.75cm},
				clip=false,
				]
				\addplot[blue] {0.5*exp(50*x)*cos(360*50*x + 90)} ;
			\end{axis}
		\end{tikzpicture}
		} % end scalebox
		
		%3T/2
			\scalebox{0.7}{
		\begin{tikzpicture}
			\begin{axis}[domain=0:60e-3, 
				samples=200, 
				xmin=0, 
				xmax=65e-3, 
				ymax=10,
				ymin=-10,
				title={Reverse Wave, $t = \frac{T}{2}$},
				scaled x ticks=base 10:3,
				xtick scale label code/.code={},
				minor x tick num = 1,
				minor y tick num = 1,
				yticklabel={\empty},
				xticklabel={\empty},
				xlabel=$z$,
				ylabel={$V(z,t)$},
				axis y line = middle,
				axis x line = middle,
				x label style={xshift=0.5cm, yshift=-0.5cm},
				y label style={xshift=-0.5cm, yshift=0.75cm},
				clip=false,
				]
				\addplot[blue] {0.5*exp(50*x)*cos(360*50*x + 180)} ;
			\end{axis}
		\end{tikzpicture}
		} % end scalebox
		
		%3T/4
			\scalebox{0.7}{
		\begin{tikzpicture}
			\begin{axis}[domain=0:60e-3, 
				samples=200, 
				xmin=0, 
				xmax=65e-3, 
				ymax=10,
				ymin=-10,
				title={Reverse Wave, $t = \frac{3}{4}$},
				scaled x ticks=base 10:3,
				xtick scale label code/.code={},
				minor x tick num = 1,
				minor y tick num = 1,
				yticklabel={\empty},
				xticklabel={\empty},
				xlabel=$z$,
				ylabel={$V(z,t)$},
				axis y line = middle,
				axis x line = middle,
				x label style={xshift=0.5cm, yshift=-0.5cm},
				y label style={xshift=-0.5cm, yshift=0.75cm},
				clip=false,
				]
				\addplot[blue] {0.5*exp(50*x)*cos(360*50*x + 270)} ;
			\end{axis}
		\end{tikzpicture}
		} % end scalebox
\end{solution}

\newpage
%%%%%%%%%%%%%%%%%%%%%%%%%%%%%%%%%%%%%%%%%
% 3.2
%%%%%%%%%%%%%%%%%%%%%%%%%%%%%%%%%%%%%%%%%
\fullwidth{\textbf{Learning Objective 3.2: }I can explain the physical meaning of and calculate the following quantities: propagation constant, phase constant, attenuation constant, characteristic impedance, VSWR, reflection coefficient, and transmission coefficient.}

% Wentworth, 2007, example 2.2, pg 46-47
\question The specifications for RG-214 are as follows:
	
		\begin{itemize}
			
			\item 1 mm \textbf{diameter} copper inner conductor 
			\item 3 mm \textbf{diameter} outer  conductor 
			\item Teflon dielectric insulator with $\varepsilon_r = 2.1$ and $\mu_r = 1$
			\item Conductivity of copper = $5.8 \ee{7}$

		\end{itemize}

\begin{parts}


	\begin{minipage}[t]{0.6\textwidth}
		\part Assuming the transmisison line is lossless, determine the distributed parameters, $R^{\prime}, L^{\prime}, G^{\prime}, C^{\prime}$
		
			\begin{solution}[4in]
			\eq{
				&R^{\prime} = 0 \\
				&L^{\prime} = \dfrac{\mu_0}{2\pi}\text{ln}\lp \dfrac{b}{a} \rp 
					= \dfrac{4\pi\ee{-7}}{2\pi}\text{ln}\lp \dfrac{1.5\text{mm}}{0.5\text{mm}} \rp
					= 2.197\ee{-7} \text{H/m}\\
				&G^{\prime} = 0 \\
				&C^{\prime} = \dfrac{2\pi\varepsilon}{\text{ln}\lp \dfrac{b}{a} \rp } 
					= \dfrac{\lp 2\pi \rp\lp 2.1 \rp \lp 8.854\ee{-12} \rp }{\text{ln}\lp\dfrac{1.5\text{mm}}{0.5\text{mm}}\rp} 
					= 1.063\ee{-10}\text{F/m}
				}
			\end{solution}
		
	\end{minipage}%
	\begin{minipage}[t]{0.3\textwidth}

			\vspace{12pt}
			
			$R^{\prime}$ = \ansbox[1.5in]{0}
	
			\vspace{6pt}
			 
			$L^{\prime}$ = \ansbox[1.5in]{v} 

			\vspace{6pt}

			$G^{\prime}$ = \ansbox[1.5in]{0} 

			\vspace{6pt}				
			
			$C^{\prime}$ = \ansbox[1.5in]{$1.063\ee{-10}\text{F/m}$} 
	\end{minipage}
	
	
	\begin{minipage}[t]{0.6\textwidth}
		\part Determine the propagation constant, attenuation constant, and phase constant. 
		
			\begin{solution}
			\eq{
				\gamma &= j\omega \sqrt{L^{\prime}C^{\prime}} 
						= j\omega \sqrt{\lp 2.197\ee{-7} \text{H/m}\rp \lp 1.063\ee{-10}\text{F/m} \rp}\\
						&= j\omega\lp 4.833\ee{-9}\rp \\
				\alpha &= 0 \\
				\beta &= \omega\lp 4.833\ee{-9}\rp 
			}
			\end{solution}
	\end{minipage}%
	\begin{minipage}[t]{0.3\textwidth}

			\vspace{12pt}
			
			$\gamma$ = \ansbox[1.5in]{$j\omega\lp 4.833\ee{-9}\rp $}
			
			\vspace{6pt}				
			 
			$\alpha$ = \ansbox[1.5in]{$0$} 
			
			\vspace{6pt}				

			$\beta$ = \ansbox[1.5in]{$\omega\lp 4.833\ee{-9}\rp $} 
				
	\end{minipage}
	
	\newpage
	\begin{minipage}[t]{0.6\textwidth}

		\part What is the characteristic impedance of this transmission line?
		
		\begin{solution}[4in]
		\eq{
			Z_0 &= \sqrt{\dfrac{L^{\prime}}{C^{\prime}}} 
				= \sqrt{\dfrac{2.197\ee{-7} \text{H/m}}{ 1.063\ee{-10}\text{F/m} }}\\
				&= 45.46\ohms
		}
		\end{solution}
	\end{minipage}%
	\begin{minipage}[t]{0.3\textwidth}
		\vspace{12pt}
		\myz{0} = \ansbox[1.5in]{$45.46$\ohms}
	\end{minipage}

	\begin{minipage}[t]{0.6\textwidth}

		\part What is the propagation velocity at 100 MHz?
		
		\begin{solution}
		\eq{
			u_p = \dfrac{\omega}{\beta} = \dfrac{c}{\sqrt{2.1}} =  0.69c = 2.07\ee{8}\text{m/s}
		}
		\end{solution}
	\end{minipage}%
	\begin{minipage}[t]{0.3\textwidth}
		\vspace{12pt}
		$u_{\text{p}}$ = \ansbox[1.5in]{0.69c}
	\end{minipage}
		


\end{parts}

\newpage
%%%%%%%%%%%%%%%%%%%%%%%%%%%%%%%%%%%%%%%%%
% 3.3
%%%%%%%%%%%%%%%%%%%%%%%%%%%%%%%%%%%%%%%%%
% Wentworth example 2.3, pg 56-57
\fullwidth{\textbf{Learning Objective 3.3: } I can use the transmission line equation to calculate the voltage/current at an arbitrary point on a specified transmission line.}

\question Consider the following \textbf{lossless} transmission line -- we want to find the voltage across the 25\ohms load. 

\scalebox{0.75}{
		\begin{circuitikz}
		\draw (0,0) to [sV, a^={$10\cosp{\omega t + 30\mydeg}$}] ++(0,3) 
				to[twoport, t=100\ohms] ++(2,0)
				to[tline,o-o] ++(3.5,0)
				to[short] ++(0.5,0)
				to[twoport,t={25\ohms}]  ++(0,-3)
				to[short] ++(-0.5,0)
				to[tline,o-o] ++(-3.5,0)
				to[short] ++(-2,0);
	%	\draw[<->] (0.5,0.25) -- node[above] {$\ell$} ++(3,0); 
		\node[below] at (3,1.75) {$\myz{0}\equals 50\ohms$};
%		\node[below] at (0.75,0) {$-\ell$};
%		\node[below] at (4.5,0) {$0$};
		\draw[->] (1.75,-0.5) --++(0,1) --++ (0.5,0);
		\node[below] at (2.25,-0.5) {\myz{in}};	
		\draw (2,3.5) --++ (3.5,0) node[midway, above] {$\ell = \lambda/4$};
		\draw (2,3.35) --++ (0,0.3);
		\draw (5.5,3.35) --++ (0,0.3);
		\end{circuitikz}
		} % end scalebox

\begin{parts}

	\begin{minipage}[t]{0.6\textwidth}

		\part Determine the load reflection coefficient
		
			\begin{solution}[2in]
			\eq{
				\Gamma_L & = \dfrac{\myz{L} - \myz{0} } {\myz{L}+\myz{}} = \dfrac{25-50}{25+50} = -\dfrac{1}{3}
			}
			\end{solution}
	\end{minipage}%
	\begin{minipage}[t]{0.3\textwidth}
		\vspace{12pt}
		$\Gamma_{\text{L}} $ = \ansbox[1.5in]{$-\dfrac{1}{3}$}
	\end{minipage}

	\begin{minipage}[t]{0.6\textwidth}
		\part Determine the input impedance (\myz{in}) for $\ell = \frac{\lambda}{4}$.
		
			\begin{solution}[2in]
			\eq{
				\myz{in} &= \myz{0}\dfrac{\myz{L} + j\myz{0}\text{tan}\lp \beta\ell \rp } {\myz{0} + j\myz{L}\text{tan}\lp \beta \ell \rp } 
				= \myz{0}\dfrac{\myz{L} + j\myz{0}\text{tan}\lp \frac{2\pi}{\lambda}\mycdot\frac{\lambda}{4} \rp } {\myz{0} + j\myz{L}\text{tan}\lp\frac{2\pi}{\lambda}\mycdot\frac{\lambda}{4}  \rp } \\
				& = \myz{0} \lp \dfrac{\myz{0}}{\myz{L}} \rp = \dfrac{\myz{0}^2}{\myz{L}} \\
				\implies \myz{in} &= \dfrac{50^2}{25} = 100\ohms
			}
			\end{solution}
	\end{minipage}%
	\begin{minipage}[t]{0.3\textwidth}
		\vspace{12pt}
		
		\myz{in} = \ansbox[1.5in]{100\ohms}
	\end{minipage}
	
	\newpage
	\begin{minipage}[t]{0.6\textwidth}
		\part Using a voltage divider (between the source impedance and the input impedance to the transmission line), determine the input voltage at $z=-\frac{\lambda}{4}$ (this will be $\widetilde{V}_{\text{in}}$ ). 
		
			\begin{solution}[3in]
			\eq{
				\widetilde{V}_{\text{in}} &= 10\angle{30\mydeg} \lp \dfrac{100\ohms}{100\ohms + 100\ohms} \rp = 5\angle{30\mydeg}
			}
			\end{solution}
	\end{minipage}%
	\begin{minipage}[t]{0.3\textwidth}
		\vspace{12pt}
		
		\myv{in} = \ansbox[1.5in]{$5\angle{30\mydeg}$}
	\end{minipage}

	\begin{minipage}[t]{0.6\textwidth}
		\part Recognizing $\widetilde{V}_{\text{in}} = \widetilde{V}(z=-\frac{\lambda}{4}) $, calculate $V_0^+$.
		
			\begin{solution}
			\eq{
				\widetilde{V} \lp z \equals -\lambda/4 \rp 
					&= 5\angle{30\mydeg} 
					= V_0^+ \lp e^{j \frac{\pi}{2}} - \dfrac{1}{3} e^{-j\frac{\pi}{2}} \rp \\
					&= V_0^+ \lp 1 \angle{90\mydeg} - \dfrac{1}{3} \angle{-90\mydeg} \rp \\
					\implies V_0^+ &= \dfrac{5\angle{30\mydeg}}{  1 \angle{90\mydeg} - \dfrac{1}{3} \angle{-90\mydeg} } = 3.75\angle{-60\mydeg}
			}
			\end{solution}
	\end{minipage}%
	\begin{minipage}[t]{0.3\textwidth}
		\vspace{12pt}
		
		$V_0^+$ = \ansbox[1.5in]{$3.75\angle{-60\mydeg}$V}
	\end{minipage}
	

	
	\newpage
	\begin{minipage}[t]{0.6\textwidth}
		\part Using these results and $\widetilde{V}(z=0) = \widetilde{V}_{\text{L}}$, calculate the voltage at the load and write the final, time-dependent equation for the voltage at the load. 
	
			\begin{solution}[3in]
			\eq{
				\widetilde{V} \lp z \equals 0 \rp 
					&= V_0^+ \lp e^{j \beta z} + \Gamma_{\text{L}} e^{-j \beta z} \rp \\
					&= 3.75\angle{-60\mydeg} \lp e^{j (0) } - \dfrac{1}{3} e^{-j (0)} \rp \\
					&=  2.513\angle{-60\mydeg}\text{ V}
			}
			\end{solution}
	\end{minipage}%
	\begin{minipage}[t]{0.3\textwidth}
		\vspace{12pt}
		
		\myv{L} = \ansbox[1.5in]{$2.513\angle{-60\mydeg}$V}
	\end{minipage}
	

	\begin{minipage}[t]{0.6\textwidth}
		\part What is the VSWR of this system? 
	
		\begin{solution}
		\eq{
			\text{VSWR} &= \dfrac{1 + \dfrac{1}{3}}{1-\dfrac{1}{3}} = 2
		}
		\end{solution}
	\end{minipage}
	\begin{minipage}[t]{0.3\textwidth}
		\vspace{12pt}
		
		VSWR = \ansbox[1.5in]{2}
	\end{minipage}		

	
\end{parts}

\newpage
%%%%%%%%%%%%%%%%%%%%%%%%%%%%%%%%%%%%%%%%%
% 3.4
%%%%%%%%%%%%%%%%%%%%%%%%%%%%%%%%%%%%%%%%%
\fullwidth{\textbf{Learning Objective 3.4: } I can use the transmission line equation to match a load to a source using a length of transmission line.}

% Wentworth 2.36 HW, pg 111

\question A matching network consists of length of 50\ohms transmission line in series with a \textbf{capacitor}. Determine the length (in wavelengths) of the transmission line and the value of the capacitor needed at 1GHz to properly match a $10-j35 \ohms$ load to a 50\ohms source. 

\begin{parts}
	
	\begin{minipage}[t]{0.6\textwidth}
		\part Normalize the impedances
	\end{minipage}%
	\begin{minipage}[t]{0.3\textwidth}
%			\vspace{12pt}
			
				\myz{s} = \ansbox[1.5in]{1\ohms}
				
				\vspace{6pt}
				
				\myz{L} = \ansbox[1.5in]{0.2 - j0.7\ohms} 

	\end{minipage}
	
	\begin{minipage}[t]{0.6\textwidth}
		\part Calculate the reflection coefficient at the load.
		
			\begin{solution}[2in]
			\eq{
				\Gamma_{\text{L}} &= \dfrac{0.2 - j0.7 -1}{0.2 - j0.7 +1 } 
					= 0.765\angle{-108.6\mydeg} 
			}
			\end{solution}
			
	\end{minipage}%
	\begin{minipage}[t]{0.3\textwidth}
		\vspace{12pt}
				$\Gamma_{L}$ = \ansbox[1.5in]{$0.765\angle{-108.6\mydeg}$} 
	\end{minipage}

	\begin{minipage}[t]{0.6\textwidth}
		\part Determine the input impedance such that the real part matches the source impedance.
		
			\begin{solution}
			\eq{
				| \Gamma_{\text{L}} | &= \dfrac{| z_{\text{in}} - 1 |}{z_{\text{in}} +1 |} \\
				0.765 & = \dfrac{ | 1 + jx -1|}{|1 + jx +1 |} \\
				0.765 &= \dfrac{ \sqrt{x^2}} { \sqrt{2^2 + x^2}} \\
				0.765^2 &= \dfrac{x^2}{4 + x^2} \\
				0.585*4 & = x^2 - 0.585x^2 \\
				\sqrt{\dfrac{2.34}{0.415}} & = x \\
				\pm 2.375 &= x
			}	
			\end{solution}
			
	\end{minipage}%
	\begin{minipage}[t]{0.3\textwidth}
		\vspace{12pt}
				\myz{in} = \ansbox[1.5in]{$1+j2.375$} 
	\end{minipage}
	
	\newpage
	\begin{minipage}[t]{0.6\textwidth}
		\part Determine the value of the capacitor that is required to compensate for the reactance of \myz{in} at 1GHz.
			\begin{solution}[3in]
			\eq{
				\dfrac{1}{j\omega C} &= -j2.375\myz{0} \\
				\implies C &= \dfrac{1}{(2\pi)(1\ee{9}\text{Hz})(2.375)*(50\ohms)} \\
					&= 1.34\text{pF}
			}
			\end{solution}
			
	\end{minipage}%
	\begin{minipage}[t]{0.3\textwidth}
		\vspace{12pt}
			C = \ansbox[1.5in]{$1.34$pF} 
	\end{minipage}
	
	
	\part Verify the length of the transmission line is $\ell = 0.294\lambda$.
	
	\begin{solution}
		\eq{
			z_{\text{in}} &=  \dfrac{z_{\text{L}} + j\text{tan}\lp \beta \ell \rp } {1 + jz_{\text{L}}\text{tan}\lp \beta \ell \rp } \\
				& = \dfrac{0.2-j0.7 + j\text{tan}\lp 2\pi\mycdot 0.294 \rp }
							{ 1 + j\lp 0.2-j0.7 \rp \text{tan}\lp 2\pi\mycdot 0.294 \rp } \\
				z_{\text{in}} & = 1.013 + j2.392\longrightarrow\text{ close enough} 
		}
	\end{solution}
\end{parts}

\end{questions}
